\subsection{Strict higher dimensions and the audit of 317}
\label{sec:F3}
\providecommand{\sSet}{\mathrm{sSet}}
\providecommand{\nCat}{n\mathrm{Cat}}
\providecommand{\Sd}{\mathrm{Sd}}
\providecommand{\Ex}{\mathrm{Ex}}
We write $\mathcal M_n$ for the structure asserted in 317,
$L_n=c_n\Sd^2$, $R_n=\Ex^2N_n$, and $B_nX=|N_nX|$.
Here $n=\infty$ means strict globular $\omega$-categories. The symbols
$W_n$ in this subsection mean the weak equivalences of $\mathcal M_n$;
they do not mean the truncation localizers also denoted by $W_n$ in
the historical discussion.

\begin{claim}{F.3a}{Proved}{Audit of the higher Thomason transfer in 317}
The arguments in the supplied source of 317 establish the proper
combinatorial structure and the Quillen equivalence
$L_n\dashv R_n$ that are used below. The following audit records the
interfaces on which the transfer depends. It concerns the supplied
mathematical proof; no Lean file was supplied or checked.
\end{claim}
\begin{proof}
The strict reflection $\tau_{\leq n}$ identifies parallel $n$-cells
which bound $(n+1)$-cells and then imposes the resulting congruence.
This is the left adjoint to inclusion of strict $n$-categories in
strict $\omega$-categories. In particular it is a reflection, not
the operation of forgetting cells above dimension $n$.
Section~2 of 317 applies the directed-chain construction to
strong Steiner complexes. The exact external results needed here
are Steiner's Theorem~5.11, giving the equivalence on complexes with
loop-free unital bases, and Theorem~6.1, giving the atomic
presentation, together with Ara--Maltsiniotis Theorem~A.15 and
A.16: the Gray tensor extends the tensor on strong Steiner complexes,
preserves all small colimits in each variable, and has the two
internal homs. These statements were opened in
\cite{I_Steiner2004,I_AraMaltsiniotis2020}. The same appendix,
Proposition~A.29, says that an internal hom into an $n$-category is
an $n$-category. Thus the reflected cylinder bijection used in 317
has the required codomain in every finite dimension; it is not an
assumption that $n\mathrm{Cat}$ is closed under the unreflected
Gray tensor.

For completeness we checked the two chain calculations that turn
this machinery into relative homotopies. The positive
Alexander--Whitney diagonal on an ordered simplex is natural for
all order maps, including degeneracies; composing with a map to
the interval gives the simplicial homotopy induced by a right Gray
cylinder. For the relative prism, if $f\leq g$ agree on a sieve
$A\subset E$ and $g$ is constant outside $A$, the cross term on
a basis simplex is
\[
 h[w_0,\ldots,w_p]
   =[f(w_0),\ldots,f(w_p),g(w_p)].
\]
Repeated entries are interpreted as zero. On the part lying in $A$
this is zero in positive degree and the interval factor is sent
through its augmentation. On the remaining part, canceling the
alternating faces gives
$dh=hd+(-1)^p(G-F)$, with $F,G$ the endpoint chain maps.
This is exactly the signed tensor differential. The remaining
penultimate-face terms coincide because $g$ is constant on the
new part and $A$ is a sieve. Consequently the cylinder is stationary
on the entire higher subcategory $U_n(A)$, as required for gluing
along an arbitrary attaching functor. This checks Section~3,
including its stationarity condition, rather than only the effect
on vertices.

In Section~4 the slice functors taking the zero and one fibres of
an $n$-category over $[1]$ preserve colimits: their right adjoints
adjoin respectively a terminal object over one and an initial
object over zero. The object functor likewise preserves colimits,
using the chaotic category as its right adjoint. The initial and
final vertex segments of an oriental are full; this follows from
the endpoints of its atomic generators and the retractions clamping
vertices to the segment. These facts compute the labelled pieces
of a pushout even when its attaching functor identifies objects.
For a cover by full sieves, an oriental simplex lies in the member
containing its last vertex, since every vertex has an arrow to that
vertex. The dual statement uses the first vertex for cosieves.
Thus Street nerves of these covers are actual simplicial unions
along monomorphisms. This verifies the gluing hypothesis used in
the next two sections.

For Section~5, let $A\subset E$ be a sieve with right adjoint
$r:E\to A$. If $E-A$ has greatest element $x$, the maps equal to
the identity on $A$ and constant respectively at $x$ and $r(x)$
on the new part give the relative zigzag
$\id_E\leq g\geq k$. The prism just checked remains stationary
after pushing out against any $U_n(A)\to X$, and gives the
retraction homotopy proving couniversality. For a finite new part,
deleting two distinct maximal new elements gives the two full
pieces needed for induction and simplicial gluing. For an infinite
new part, finite subsets closed under $ir$ form a filtered family;
the attached category $X$ is retained in every member. Nerves
preserve these filtered colimits because each oriental is finitely
presented. The same induction, starting with a greatest element,
proves the unit $NP\to N_n c_nNP$ for every poset $P$.

Section~6 applies this to the flag posets of subdivided boundaries
and horns. Their inclusions are sieves, and on the cosieve generated
by an old flag the largest old subflag is the stated right adjoint.
First attaching this cosieve gives the couniversal weak equivalence.
The remaining pieces form full cosieve covers, so their nerves
give the monomorphic gluing square. It follows both that a pushed-out
subdivided horn is a weak equivalence and that a boundary attachment
preserves a weak equivalence between the objects to which it is
attached. These are the two different assertions used for transfer
and for left properness.

The presentability appendix treats strict categories as models of a
finitary essentially algebraic theory. Its presentation construction
simultaneously imposes equations and creates compositions whenever
the equations make their boundaries match. Thus it avoids the
incorrect procedure of merely quotienting a previously constructed
free category. All formation rules have finitely many premises,
also in the countable-sorted $\omega$-case. Finite presentations
therefore give a small dense set of finitely presentable objects,
and filtered colimits are computed cellwise. This supplies the
smallness and filtered-colimit assertions used in Section~7.

Finally $L_nI$-injectives and $L_nJ$-injectives are respectively
$R_n^{-1}(\mathrm{TrivFib})$ and $R_n^{-1}(\mathrm{Fib})$ by
adjunction. Relative $L_nJ$-cell maps are weak equivalences by the
horn assertion and filtered-colimit closure. The small-object
factorizations and the usual arrow-retract argument consequently
give both weak factorization systems. Right properness follows
because $R_n$ preserves pullbacks and detects weak equivalences,
and simplicial sets are right proper. Left properness follows by
the boundary-attachment assertion, transfinite composition and
retracts. The unit after $\Sd^2$ is proved by boundary induction,
using the already checked poset unit on each cell; composing with
the subdivision comparison gives the derived adjunction unit.
Since $R_n$ detects weak equivalences, this proves the Quillen
equivalence. These are the contents of 317, Sections~2--7 and
Appendix~A, used here. At $n=1$ the model structure and properness
also agree with Thomason's Theorem~4.9 and Corollary~5.5
\cite{I_Thomason1980}.
\end{proof}

\begin{claim}{F.3b}{Proved}{The correspondence for set-generated localizations}
Restriction along $\Cat\subset n\mathrm{Cat}$ gives an
inclusion-preserving bijection between the weak-equivalence classes
of left Bousfield localizations of $\mathcal M_n$ at small sets of
maps and the accessible fundamental localizers of $\Cat$.
For an accessible fundamental localizer $W$, the inverse is
\[
 W^{(n)}=\{f:i_\Delta N_n(f)\in W\},                 \tag{F.3.1}
\]
where $i_\Delta K=\Delta/K$ is the category of simplices.
Equivalently $B_n(f)$ belongs to the spatial localization
corresponding to $W$. The same assertion holds if ``at small sets''
is replaced by ``combinatorial''.
\end{claim}
\begin{proof}
We use two precise existence statements. Cisinski's
Theorem~4.2.15 identifies fundamental localizers modelled by a test
category with its test localizers and preserves accessibility.
For the test category $\Delta$, its inverse is $i_\Delta^{-1}$;
the ordinary nerve gives the forward class on $\Cat$, by
Remark~4.2.16 and the test-functor statement used there.
Theorem~1.4.3 identifies accessible simplicial localizers with the
cofibrantly generated model structures having monomorphisms as
cofibrations. Thus accessible fundamental localizers correspond
exactly to set-generated localizations of spaces. We use this
specific form of \cite{I_Cisinski2006}.

Barwick's Theorem~4.7 states that a left proper combinatorial model
category localized at a small set of homotopy classes has a left
proper combinatorial model structure, with the original
cofibrations, the fibrant local objects as fibrant objects, and the
local equivalences as weak equivalences \cite{I_Barwick2010}.
The theorem uses derived mapping spaces and does not require the
original model category to be simplicial. It applies to $\mathcal M_n$
by F.3a.

Let $S$ be a small set of simplicial maps generating the spatial
localization associated to $W$. Every simplicial set is cofibrant,
so the derived images of its maps under $L_n$ are represented by
the actual maps $L_nS$. For a fibrant $Z$ in $\mathcal M_n$,
the derived mapping adjunction gives
\[
 \Map^h_{\mathcal M_n}(L_nK,Z)
       \simeq\Map^h_{\sSet}(K,R_nZ).
\]
It follows, for each map in $S$ and each such $Z$, that $Z$ is
$L_nS$-local exactly when $R_nZ$ is $S$-local. The Quillen
equivalence in F.3a is essentially surjective on fibrant homotopy
types. Thus testing a map of $\mathcal M_n$ against every local
object is exactly testing its image in spaces against every
$S$-local object. This proves formula (F.3.1) and its existence.

Conversely, a set of maps in $\mathcal M_n$ represents a small
set of maps in spaces under the same equivalence. The mapping-space
argument in the opposite direction proves that its local-equivalence
class is the inverse image of that spatial class. On an ordinary
category, $N_n$ is exactly the ordinary nerve. The natural
last-vertex comparison $i_\Delta NA\to A$ is a weak equivalence;
its compatibility with functors shows that the restriction in
(F.3.1) is exactly $W$. In particular no choices of replacements
change this class. This proves surjectivity and injectivity.

For the last assertion, a combinatorial localization has a small
set $J'$ of generating trivial cofibrations. Its fibrant objects
are precisely the original fibrant objects local with respect to
$J'$: one direction follows from the lifting property, and the
other from the characterization of local equivalences by derived
mapping spaces and the two weak factorization systems. Equivalently,
apply the standard generating-trivial-cofibration proof of
\cite[Proposition~A.3.7.4]{I_LurieHTT}, after replacing the derived
mapping spaces by cosimplicial resolutions. This proof uses no
simplicial enrichment of $\mathcal M_n$. Hence this localization
has the same weak equivalences as localization at $J'$.
\end{proof}

\begin{claim}{F.3c}{Refuted as an unrestricted ZFC assertion}{A nonaccessible model localization}
Suppose the indexing universe $\mathcal U$ contains no measurable
cardinal. In every strict dimension $1\leq n\leq\infty$, there
is a left Bousfield localization of $\mathcal M_n$ whose
weak-equivalence class restricts to a nonaccessible fundamental
localizer of $\Cat$. Thus the unrestricted wording in Part F.3,
which does not require the localized model structure to be
combinatorial, has a counterexample in this ZFC-compatible setting.
Under Vop\v{e}nka's principle for the relevant $\mathcal U$-classes,
the unrestricted correspondence is a bijection. The unconditional
size-correct correspondence is F.3b.
\end{claim}
\begin{proof}
Put $A_\lambda=\ZZ^\lambda/\ZZ^{<\lambda}$ for infinite
$\mathcal U$-small cardinals $\lambda$, where $\ZZ^{<\lambda}$
consists of the vectors whose support has cardinality less than
$\lambda$, and let $\mathcal A$
be their class. A group is reduced when it admits no nontrivial
homomorphism from any $A_\lambda$. Reduced groups are closed
under subgroups, products and extensions. For the extension
assertion, a map $A_\lambda\to G$ into an extension of two
reduced groups first has trivial image in its quotient and then
trivial image in its kernel.

For a $\mathcal U$-small group $G$ let
\[
 R(G)=\bigcap\{\ker(a):a:G\to H,\ H\text{ reduced}\},
 \qquad P G=G/R(G).
\]
This construction is small: it suffices to use quotient images
of $G$, and their normal kernels form a subset of its power set.
The diagonal embedding into the product of these reduced quotients
shows that $PG$ is reduced. Every map from $G$ to a reduced group
factors uniquely through $PG$, giving a functorial reflector.
It applies to groupoids by quotienting each isotropy group by its
characteristic radical; conjugation carries these radicals to
one another, so the arrow torsors and composition descend.

We need the following additional fact to produce an actual model
structure, rather than merely a homotopy reflector:
\[
                 R(R(G))=R(G).                              \tag{F.3.3}
\]
Indeed $R(R(G))$ is characteristic in the normal subgroup $R(G)$,
and hence normal in $G$. The quotient $G/R(R(G))$ is an extension
of the reduced groups $P(R(G))$ and $PG$, so is reduced. The
universal property of $R(G)$ forces $R(G)\subseteq R(R(G))$.
The reverse inclusion is immediate. Thus $P(R(G))=1$.

Let $\Pi^{(n)}_1X$ be the fundamental groupoid obtained by
reflecting a strict $n$-category to an ordinary category and
inverting all arrows, and regard
\[
                   Q_nX=P\Pi^{(n)}_1X
\]
as a strict $n$-category concentrated in invertible dimensions zero
and one. There is a natural strict functor $X\to Q_nX$.
The zero, one and two simplices of the Street nerve give exactly
the same presentation of the fundamental groupoid: its edges are
the $1$-cells, and a triangular $2$-simplex imposes the relation
between its two composite boundary paths. Consequently
\[
 \Pi_1(B_nX)\simeq\Pi^{(n)}_1X,
 \qquad B_nQ_nX\simeq N P\Pi_1(B_nX).
\]
The functor $Q_n$ preserves $W_n$, and its unit is idempotent
up to the canonical isomorphism. Its local homotopy types are
precisely the $1$-types with reduced fundamental groups.

We apply the following exact form of the Bousfield--Friedlander
criterion, Stanculescu's Theorem~1.1 \cite{I_Stanculescu2008}.
For a model category and a coaugmented endofunctor $Q$, assume
that $Q$ preserves weak equivalences, that $Q\eta_X$ is a weak
equivalence and $\eta_{QX}$ is monic in the homotopy category,
and that $Q$-equivalences are preserved by pullback along original
fibrations between original fibrant $Q$-local objects. Then there
is a left Bousfield model localization with the original
cofibrations and the $Q$-equivalences as weak equivalences.
No small generating family is assumed. We opened the corrected
primary version, including its factorization proof.

The first two hypotheses were just verified. For the third,
an original fibration between fibrant objects has strict pullbacks
which $R_n=\Ex^2N_n$ identifies with homotopy pullbacks of
spaces. It therefore suffices to check the asserted property for
a homotopy pullback of a $Q$-equivalence $X\to B$ along
$E\to B$, where $E,B$ are local $1$-types. Work on one component
of $E$. The map on fundamental groups $g:G\to H$ induced by
$X\to B$ is a surjection with kernel $R(G)$: since $H$ is
reduced it factors as $G\to PG\xrightarrow{\sim}H$.
Surjectivity makes the homotopy pullback connected over this
component. Since $\pi_2B=0$, its fundamental group is
\[
             G\times_H K\longrightarrow K,
             \qquad K=\pi_1E.
\]
This projection is surjective with kernel $R(G)$. Every map
from this pullback group to a reduced group kills its kernel,
by (F.3.3), and hence factors uniquely through $K$. Since $K$
is reduced, applying $P$ to the projection is an isomorphism.
The same argument on each component proves the required
$Q$-equivalence. The criterion now constructs the promised
model localization of $\mathcal M_n$. In particular its existence
does not rely on an unproved transfer of a proper-class reflector
to a model structure.

It remains to verify that the restriction $W_Q$ to $\Cat$ is
a fundamental localizer and is nonaccessible. Mapping into all
reduced groupoids detects $P\Pi_1$-equivalences, by the reflector
universal property and Yoneda for groupoids. These are exactly
the derived mapping-space tests against the local $1$-types.
For any fixed target space such tests define a fundamental
localizer: the two elementary axioms are immediate, and relative
Theorem A follows from
$BA\simeq\hocolim_{s\in S}B(A/s)$ and the corresponding
formula for $B$, since mapping into the fixed target turns this
homotopy colimit into a homotopy limit. Intersecting these classes
over the local targets proves that $W_Q$ is fundamental.

Proposition~6.2 of Casacuberta--Scevenels--Smith states that the
absence of measurable cardinals is equivalent to $\ZZ$ being
$\mathcal A$-reduced \cite{I_CSS2005}. The statement relativizes
to $\mathcal U$: a nonzero map $A_\lambda\to\ZZ$ would produce
a measurable cardinal at most $\lambda$. Therefore every
$W_Q$-map induces isomorphisms on integral cohomology in degrees
zero and one. Degree zero follows from the bijection of
components; degree one follows, on each component, from
$\Hom(PG,\ZZ)=\Hom(G,\ZZ)$.

Suppose a $\mathcal U$-small set $S$ generated $W_Q$. Choose
an uncountable regular $\kappa\in\mathcal U$ larger than the
total number of simplices in the sources and targets of $NS$.
The functor $T_\kappa(B)=B^\kappa/B^{<\kappa}$ on abelian
groups is exact: both products and the union of the bounded-support
subproducts preserve short exact sequences, and the quotient
does so by the snake lemma. If a free chain complex $C_*$ has
fewer than $\kappa$ basis elements in total, the interchange of
its two coordinate sets gives
\[
 \Hom(C_*,A_\kappa)
       \cong T_\kappa\Hom(C_*,\ZZ).
\]
For the kernel assertion here, the union of the supports of
fewer than $\kappa$ basis values still has size less than
$\kappa$; this is exactly where regularity is used. Thus
$H^j(C_*;A_\kappa)\cong T_\kappa H^j(C_*;\ZZ)$.
Applying this in degrees zero and one to $NS$ shows that
$Z=K(A_\kappa,1)$ is $S$-local. Indeed its mapping spaces have
components $H^1(-;A_\kappa)$, fundamental groups
$H^0(-;A_\kappa)$, and no higher groups; all their components
are translates.

The fixed-target test against $Z$ is a fundamental localizer
containing $S$, so it must contain $W_Q$. But the one-object
groupoid $A_\kappa\to *$ belongs to $W_Q$, because
$P A_\kappa=1$. Its mapping-space test against $Z$ fails:
the identity of $K(A_\kappa,1)$ and the constant map induce
different maps on $\pi_1$, since $A_\kappa\ne0$. This
contradiction proves nonaccessibility. It also gives a direct
bounded proof of the required form of the obstruction in
\cite[Theorems~6.3--6.4]{I_CSS2005}.

Finally an arbitrary model localization has its original-fibrant
local objects as a full homotopy subcategory, and its weak
equivalences are detected by derived mapping spaces into them.
This follows using the same cofibrant replacements and original
cosimplicial resolutions: the cofibrations and trivial fibrations
have not changed, and an original resolution remains a resolution
in the localized structure. After the equivalence with spaces in
F.3a, restriction to $\Cat$ is therefore a fundamental localizer
by the preceding relative-Theorem-A argument. Under the class
form of Vop\v{e}nka used in Part H, Cisinski's Scholie~4.2.17
makes it accessible. F.3b then supplies its set-generated higher
model, and the mapping-space tests make the weak-equivalence
classes equal. This proves the positive assertion under that
principle. The two additional set-theoretic hypotheses are being
distinguished explicitly; automatic accessibility in ZFC has not
been inserted into the statement of 317.
\end{proof}

\begin{claim}{F.3d}{Proved}{Reversibility and spectral classes transport exactly}
A full sieve of a strict $n$-category is the full subcategory on a
set of objects closed under taking the source of a $1$-cell whose
target is in that set. A two-sieve Mayer--Vietoris square consists
of a cover $X_3=X_1\cup X_2$ of this kind and $X_0=X_1\cap X_2$.
For a map of two such squares, higher Loc~$(6\,b)$ is the assertion
\[
 f_1,f_2,f_3\in W^{(n)}\quad\Longrightarrow\quad f_0\in W^{(n)}.
\]
For the localizers in F.3b this holds exactly when $W$ satisfies
ordinary Loc~$(6\,b)$. For a localizing subcategory $\mathcal L$
of spectra put
\[
 R_n(\mathcal L)=
 \{f:\Sigma^\infty C(B_nf)\in\mathcal L\}.           \tag{F.3.2}
\]
Then $W^{(n)}=R_n(\mathcal L)$ if and only if
$W=R(\mathcal L)$. Therefore the accessible classification question,
including the negative answer of Part B, has the same answer in
every strict dimension.
\end{claim}
\begin{proof}
The last-vertex argument in the audit shows that applying $N_n$
to a two-sieve cover gives a union of simplicial subsets.
Applying $i_\Delta$ gives a cover by full ordinary sieves: every
face or degeneracy of a simplex lying in one member still lies in
that member. Thus a higher Mayer--Vietoris cube becomes an ordinary
Mayer--Vietoris cube, with its four comparison functors exactly
$i_\Delta N_nf_i$. Formula (F.3.1) proves the implication from
ordinary Loc~$(6\,b)$ to the higher version. Conversely ordinary
categories and their full sieves are special cases of the higher
ones, so restriction proves the reverse implication. This argument
also proves the correspondence for the forward gluing axiom and
for the reverse single-square axiom, without mistaking an arbitrary
algebraic pushout for a homotopy pushout.

The category-of-simplices comparison gives a natural weak
equivalence $B(i_\Delta N_nX)\simeq B_nX$. It is natural in the
specified map $f$, hence induces an equivalence of its homotopy
cofibers, including the convention for empty sources. Formula
(F.3.2) is consequently exactly the inverse-nerve class of
$R(\mathcal L)$. Restriction gives necessity, and (F.3.1) gives
sufficiency. In particular applying F.3b to the accessible
counterexample of Part B produces the corresponding counterexample
in every dimension. No assertion that arbitrary $R_n(\mathcal L)$
is set-generated is needed for this equivalence.
\end{proof}

\begin{claim}{F.3e}{Proved}{What properness gives for the fibration axioms}
For every $1\leq n\leq\infty$, right properness in 317 implies
Loc~$(7\,F)$ for $W_n$, with strict pullbacks and with
\[
 \Fib_{W_n}=
 \{p:\text{every base change along }p\text{ preserves }W_n\}.
\]
For $n=1$, $W_1=W_\infty$ also satisfies Grothendieck's stronger
Loc~$(7)$ and Loc~$(7\,D)$. Their proof uses the simplex test for
the ordinary nerve, in addition to properness. In higher dimensions,
the direct conclusion of the proper model structure is the same
Loc~$(7\,F)$; 317 supplies no higher definition of the Cat-fibration
or perfect-functor classes occurring in the literal other two axioms.
\end{claim}
\begin{proof}
Every model fibration lies in $\Fib_{W_n}$ by right properness.
Factor any strict $n$-functor as a model trivial cofibration
followed by a model fibration. This gives a $W_n$-map followed
by a member of $\Fib_{W_n}$, the existence assertion in
Loc~$(7\,F)$. Composition and strict base change preserve this
class by its definition. Every map $X\to *$ is in it, because
the Street nerve preserves products and products with a fixed
simplicial set preserve weak equivalences. Isomorphisms are in it;
pullbacks of its maps and of its weak equivalences exist and have
the required properties. Finally factoring the diagonal provides
the path-object axiom. These verifications give the stated Brown
category of fibrations, including the terminal and empty cases.

For ordinary categories, let $p:E\to B$ satisfy Grothendieck's
weak-fibration test (Ff), XVI ms.~27: pullback preserves weak
equivalences between categorically contractible bases. Every
ordinal category $[m]$ is contractible. Hence, for every simplex
$[n]\to B$ and order map $[m]\to[n]$, (Ff) gives a weak
equivalence
\[
 E\times_B[m]\longrightarrow E\times_B[n].
\]
Since $N[m]=\Delta[m]$ and $N$ preserves pullbacks, these are
exactly the simplex pullback tests in Rezk's Theorem~4.1(4),
specialized to the topos of sets \cite{I_Rezk1998}. That theorem
makes $Np$ sharp, which by definition means that all its base
changes preserve weak equivalences. Taking nerves again shows
$p\in\Fib_{W_\infty}$. Thus
$\mathrm{Fib}^{\mathrm{weak}}_{W_\infty}\subseteq
\Fib_{W_\infty}$, the strongest weak-fibration formulation of
Loc~$(7)$ at XVI ms.~32. The Cat-fibrations with homotopically
constant transports and the perfect functors are subclasses of
this weak-fibration class, so the displayed inclusion gives
Loc~$(7)$ in its Cat-fibration formulation and Loc~$(7\,D)$.
For the former subclass one can check (Ff) directly by the
Grothendieck construction of the fibre diagram; for the latter
it is the smooth-and-proper base-change test in the definition of
perfect functors. Thus no claim that an arbitrary model fibration
is a Grothendieck Cat-fibration is being made.

For $n>1$, an oriental is not the ordinal category $[m]$, and
$N_n\mathcal O_m$ is not $\Delta[m]$. Thus repeating the preceding
proof with orientals would replace the actual simplex pullbacks
by different ones and would require a new comparison theorem.
Moreover the manuscript's classes $\Catfib^W$ and $\Parf_W$
are classes of ordinary functors, not definitions on strict
$n$-functors. If they are transported by the ordinary category-of-
simplices construction, the corresponding nerve-based statements
follow from the preceding paragraph by definition; that does not
identify an intrinsic higher Cat-fibration class. The exact
dimension-independent implication of 317 itself is therefore the
proved Loc~$(7\,F)$. Left properness gives pushout invariance
along model cofibrations, and supplies no additional assertion
about those stronger fibration classes.
\end{proof}
